% Source fingerprint: a7c7f4ff5bc13a4c525ed53af52b7a2c445754c15aa2a37567b3311dae02b8a3
% T01: raw TeX cells preserved; no numeric highlighting.
\begin{table}[htbp]
\centering\small\renewcommand{\arraystretch}{1.18}\setlength{\tabcolsep}{4pt}
\caption[集合系闭性与生成工具]{集合系闭性与生成工具。各行只列定义性要求；生成工具的前提不同，不能用任意生成集合族代替定理中的代数或交封闭族。}
\label{b1:tab:set-system-closures}
\begin{tabularx}{\linewidth}{@{}>{\raggedright\arraybackslash}p{3.1cm}>{\raggedright\arraybackslash}p{4.7cm}>{\raggedright\arraybackslash}X@{}}
\toprule
集合系 & 定义性要求 & 生成与推广工具 \\
\midrule
集合代数 & 补集、有限并 & 有限集合运算；从代数出发用单调类定理得到生成的 \(\sigma\)-代数 \\
\(\sigma\)-代数 & 补集、可数并 & 测度和可测映射的完整定义域；生成对象按最小性构造 \\
\(\pi\)-系统 & 有限交 & 包含它的 Dynkin 系统包含其生成的 \(\sigma\)-代数，是唯一性推广的入口 \\
Dynkin \(\lambda\)-系统 & 全集、嵌套差、不交可数并 & 须与 \(\pi\)-系统条件配合，不能单凭这些闭性当作 \(\sigma\)-代数 \\
单调类 & 递增并、递减交 & 包含一个集合代数时，包含该代数生成的全部可测集 \\
\bottomrule
\end{tabularx}
\end{table}
